\chapter{Implementation}
\label{ch:implementation}

\section{Ingestion}
\label{sec:impl-ingestion}

\subsection{Partitioning Is a Correctness Decision}
\label{sec:impl-partitioning}

The topic \code{fleet.telemetry} has 6 partitions, \textbf{keyed by
\code{vehicle\_id}}. This is not about throughput. Kafka guarantees
ordering only \emph{within} a partition~\cite{kreps2011kafka}, and the
speed layer's per-vehicle state machine --- which tracks when each vehicle
last became idle --- assumes it observes one vehicle's events in the order
they occurred. Keying by \code{vehicle\_id} provides exactly that
guarantee. A round-robin key would silently corrupt every idle
measurement, and would do so without producing a single error.

The dead-letter topic \code{fleet.telemetry.dlq} has 1 partition: it
needs no ordering and carries a trickle of traffic, so one partition
keeps every rejected event in a single readable stream.

\subsection{Deliberate Fault Injection}
\label{sec:impl-faults}

The simulator corrupts its own output so that error handling is
\emph{demonstrable} rather than asserted. Table~\ref{tab:faults} gives
configured against measured rates over approximately 2,700 events.

\begin{table}[htbp]
\centering
\caption{Injected fault rates, configured against measured, and what
each fault demonstrates.}
\label{tab:faults}
\begin{tabularx}{\textwidth}{@{}llrX@{}}
\toprule
\textbf{Fault} & \textbf{Configured} & \textbf{Measured} &
\textbf{What it proves}\\
\midrule
Malformed JSON & 0.5\% & 0.40\% &
  The archiver preserves unparseable bytes; the speed layer
  dead-letters rather than crashing\\
Schema-invalid & 0.5\% & 0.77\% &
  \code{common/validation.py} rejects identically in both layers\\
Duplicate \code{event\_id} & 1.0\% & 1.03\% &
  Deduplication works; both copies are \emph{valid}, so nothing else
  can stop them\\
Late by 1--20 sim min & 2.0\% & 2.09\% &
  The consistency argument --- see Section~\ref{sec:consistency}\\
\bottomrule
\end{tabularx}
\end{table}

The DLQ carries the reason code \emph{and the original payload}, so
rejected events are replayable once a producer is fixed. The observed
distribution across five reason codes was \code{unparseable\_json}~58,
\code{coords\_out\_of\_bounds}~23, \code{bad\_status}~17,
\code{negative\_fare}~12 and \code{negative\_speed}~10.

Validation \textbf{rejects rather than repairs}. Coercing a malformed
field would produce a confidently wrong money figure downstream, which is
strictly worse than a visible rejection.

\subsection{The Daily Batch Source}

A second simulator drops one expense CSV per simulated day, carrying
\code{vehicle\_id}, \code{fuel\_cost}, \code{maintenance\_cost},
\code{distance\_covered} and \code{service\_flag}. Files are versioned:
a partner correcting an earlier submission drops a \code{v2} file for a
date already reconciled, which is the event the entire architecture
exists to absorb (Section~\ref{sec:recompute}).

The service-level agreement for delivery governs \textbf{first delivery
only}. An early implementation evaluated the SLA against every version,
so a corrected \code{v2} --- by definition sent after the deadline ---
marked every resubmission late. That is discussed as defect~5 in
Section~\ref{sec:defects}.

\section{The Speed Layer}
\label{sec:impl-speed}

\subsection{The Archiver Is Deliberately Boring}
\label{sec:impl-archiver}

The archiver is a \emph{separate Spark application} with its own
checkpoint and \textbf{no business logic whatsoever}. The master dataset
is the source of truth from which every past day is recomputed; if it is
corrupted, all history is wrong and nothing can recover it. Isolating it
means a bug in the speed layer's validation, windowing or state handling
\emph{cannot reach it} --- the worst such a bug can do is produce wrong
dashboards, which the next batch run overwrites.

Unparseable records are archived under \code{sim\_date=unknown} with
their raw bytes intact. An archive that silently discarded malformed
records would make ``how much bad data did we receive?'' permanently
unanswerable, which is exactly the property immutability is supposed to
buy~\cite{helland2015immutability}.

\subsection{Stateful Idle Alerts}
\label{sec:impl-idle}

``Has this vehicle been idle for 45 simulated minutes?'' cannot be
answered by a window or an aggregation. It requires memory of \emph{when}
the vehicle last stopped being idle, carried across micro-batches, and
the ability to fire when \textbf{no} new event arrives --- a silent
vehicle being precisely the case of interest.
\code{applyInPandasWithState} is the only Structured Streaming construct
providing both~\cite{spark_ss_guide}.

Four subtleties, each a defect if mishandled:

\begin{enumerate}
  \item \textbf{The idle stopwatch starts on the \emph{transition} into
        idle}, not on every idle ping. Resetting per ping means the timer
        never advances and no alert can ever fire.
  \item \textbf{Events are sorted by event time before folding}, because
        a micro-batch may interleave partitions.
  \item \textbf{A stale out-of-order event does not rewind the state.} It
        is not lost --- the archive holds it and the batch layer counts it
        --- but rewinding would corrupt the idle measurement.
  \item \textbf{\code{EventTimeTimeout}, not \code{ProcessingTimeTimeout}.}
        A processing-time timeout fires after $N$ \emph{real} seconds, so
        changing the clock compression would silently change the business
        rule.
\end{enumerate}

A known consequence, stated rather than hidden: an event-time timeout
only fires when the watermark advances, and the watermark only advances
when other events arrive~\cite{akidau2015dataflow}. Were the entire fleet
to go silent simultaneously, no timeout would fire. That case is covered
by a different control --- the \code{TelemetryNotProduced} alert, which
watches the producer directly.

\subsection{Idempotent Sinks}
\label{sec:impl-sinks}

\code{foreachBatch} provides \textbf{at-least-once} delivery. If the
driver dies after writing a micro-batch but before committing offsets,
Spark re-runs that batch and the sink observes the same rows
twice~\cite{armbrust2018structured}. Under plain inserts, every dashboard
figure would ratchet upward on each restart and never recover.

Every sink therefore upserts on a \textbf{natural key}:
\code{rt\_vehicle\_state} on \code{(vehicle\_id)}, \code{rt\_zone\_hourly}
on \code{(zone\_id, window\_start)}, \code{rt\_vehicle\_daily} on
\code{(vehicle\_id, sim\_date)}, and \code{idle\_alerts} on
\code{(vehicle\_id, opened\_at)}. The repeated write sets identical
values, making the sink idempotent --- the practical equivalent of
exactly-once semantics~\cite{kleppmann2017ddia}.

This also permits \code{update} output mode for the windowed
aggregations, which is necessary because \code{append} would withhold a
window until the watermark passed its end, meaning the \emph{current}
simulated hour --- the thing a live dashboard exists to show --- would
never appear.

\section{The Batch Layer and Orchestration}
\label{sec:impl-batch}

\subsection{The Batch Layer Trusts Nothing}
\label{sec:impl-trust}

\code{compute\_vehicle\_day} re-reads the raw archive and
\textbf{re-applies the same validation and deduplication from scratch}.
It does not read the speed layer's output at all. Critically it applies
\textbf{no watermark}: it reads a complete, closed day, so lateness is
irrelevant to it, and it deduplicates across the entire day rather than
within a watermark window. This independence is what makes the drift
measurement in Section~\ref{sec:consistency} meaningful --- the two
figures are genuinely arrived at twice.

Its outputs per vehicle-day are \code{trips}, \code{revenue},
\code{gps\_km} (haversine between consecutive pings, with physically
impossible jumps above 150\,km/h discarded), \code{online\_min},
\code{on\_trip\_min}, \code{idle\_min} and \code{utilization}.

The reconciliation joins telemetry to expenses with a \textbf{full outer
join}, deliberately. An inner join would silently drop the two findings
the business most needs: telemetry with no expense row
(\code{missing\_costs} --- we are blind to a vehicle's costs) and an
expense row with no telemetry (\code{missing\_telemetry} --- we are being
billed for a vehicle that never reported).

\subsection{The Schedule Is a Poll, Not a Cron}

Airflow's scheduler lives in real time and knows nothing about the
compressed clock. The DAG therefore runs every 2 real minutes and its
\emph{first} task decides which simulated day is actually pending;
the schedule is a poll and the date selection is data-driven. Expressing
the schedule in simulated time was tried, and produced an interval 60
times longer than intended --- a variant of defect~3 in
Section~\ref{sec:defects}.

The DAG has nine tasks: \code{find\_pending\_date},
\code{check\_expense\_file}, \code{validate\_expenses},
\code{compute\_vehicle\_day}, \code{reconcile\_profitability},
\code{load\_batch\_views}, \code{compute\_drift},
\code{generate\_report} and \code{record\_run}.

\subsection{Recomputation Is Delete-then-Insert}

\code{load\_batch\_views} deletes every row for the target
\code{sim\_date} and inserts the recomputed set inside a single
transaction. Rerunning a date therefore produces the same row count with
updated values, and a partially applied recompute is impossible. This is
the property the architecture exists to provide, and
Section~\ref{sec:recompute} demonstrates it on both a corrected
resubmission and a forced replay.

\section{Storage and Serving}
\label{sec:impl-serving}

Both layers' outputs live in one PostgreSQL instance: four \code{rt\_*}
tables written by the speed layer and four \code{batch\_*} and
bookkeeping tables written by the batch layer
(Appendix~\ref{app:schema}). Every \code{rt\_*} table has a primary key
on its natural key, which is what makes the streaming upserts of
Section~\ref{sec:impl-sinks} possible.

\subsection{The Lambda Merge Is Explicit, Never Blended}
\label{sec:impl-merge}

\code{GET /vehicles/\{id\}} returns three blocks, each labelled with its
\code{source} and its \code{as\_of}:

\begin{lstlisting}[language=,basicstyle=\ttfamily\scriptsize]
{ "live":    { "source": "speed", "as_of": "...", "status": "idle" },
  "today":   { "source": "speed", "trips": 4, "revenue": 512.25 },
  "history": [ { "source": "batch", "sim_date": "...", "net_profit": 570.0 } ] }
\end{lstlisting}

They are deliberately \textbf{not merged into a single figure}. The two
layers answer different questions with different guarantees, and a
consumer must be able to tell at a glance which numbers are safe to quote
to finance. Blending them is precisely how a dashboard estimate ends up
in a financial statement.

``Now'' is always \emph{simulated} now, derived from
\code{max(last\_event\_time)} rather than the database clock --- so if
ingestion stalls, ``now'' stops advancing and the active-vehicle count
correctly falls to zero rather than quietly lying.

The API exposes eleven documented endpoints
(Appendix~\ref{app:api}), and the consolidated daily report is served as
a rendered HTML document at \code{GET /reports/\{sim\_date\}}. Two
Grafana dashboards --- \emph{Fleet operations} and \emph{Pipeline health}
--- read the same tables.

\section{Observability}
\label{sec:impl-observability}

\subsection{What Is Measured, and Why}

Three questions drive every metric. Nothing is collected because it was
easy to collect.

\begin{table}[htbp]
\centering
\caption{The metric design. Each row answers a question that the others
cannot.}
\label{tab:metrics}
\begin{tabularx}{\textwidth}{@{}lXX@{}}
\toprule
\textbf{Question} & \textbf{Metrics} & \textbf{Why it is not obvious}\\
\midrule
Is data arriving? &
  \code{fleet\_events\_produced\_total}, \code{fleet\_faults\_injected\_total} &
  Distinguishes ``the producer is silent'' from ``the producer is gone''
  --- the alert needs both \code{increase()==0} and \code{absent()}.\\
\addlinespace
Is it being processed? &
  \code{fleet\_stream\_last\_progress\_timestamp},
  \code{\ldots\_processed\_rows\_per\_second},
  \code{\ldots\_offsets\_behind\_latest} &
  Freshness is per-query, not per-process. A container can be perfectly
  healthy while one of its five queries is frozen.\\
\addlinespace
Is it correct? &
  \code{fleet\_dlq\_events\_total}, \code{fleet\_late\_rows\_dropped\_total},
  \code{fleet\_speed\_batch\_drift\_ratio} &
  These measure \emph{known, accepted} imprecision. A pipeline that
  cannot quantify its own error cannot be trusted with money.\\
\addlinespace
Is the business healthy? &
  \code{fleet\_open\_idle\_alerts}, the profitability flags &
  Deliberately separate from pipeline health:
  \code{IdleVehiclesHigh} fires when everything is working and the fleet
  simply is not earning.\\
\bottomrule
\end{tabularx}
\end{table}

\subsection{Structured Logging Across All Stages}

Every log line is a single JSON object carrying \code{ts}, \code{level},
\code{service}, \textbf{\code{stage}}
(\code{ingestion}/\code{processing}/\code{storage}/\code{serving}/\code{orchestration}),
\code{event}, \code{msg} and \textbf{\code{sim\_time}}.

\code{sim\_time} is the field that matters here. On a compressed clock a
wall-clock timestamp tells you nothing about \emph{which simulated hour} a
line belongs to, so every line carries both. Per-event logging is
sampled 1-in-$N$: at 50 vehicles emitting every 2 seconds, an INFO line
per event would bury everything else, so the useful granularity is one
line per micro-batch carrying counts.

\subsection{Tracing --- An Honest Position}
\label{sec:tracing}

We do \textbf{not} implement distributed tracing. What we implement is
\textbf{correlation-ID propagation}: \code{event\_id} travels producer
$\rightarrow$ Kafka $\rightarrow$ archiver $\rightarrow$ DLQ
$\rightarrow$ batch dedup; \code{vehicle\_id} travels through every stage
and is the partition key; \code{run\_id} tags every Airflow task,
\code{batch\_runs}, \code{dq\_issues} and \code{batch\_vehicle\_daily};
and \code{sim\_date} identifies the lake partition, every batch table and
the report filename.

A single rejected event can therefore be traced from its DLQ reason back
to its original payload, and a single reconciled figure back to the run
and expense-file version that produced it. What this does \emph{not} give
is per-request span timing across service boundaries. For a pipeline
whose hops are Kafka offsets and Parquet partitions rather than RPC
calls, spans would add infrastructure without answering a question we
actually have. Section~\ref{sec:production} states when that calculus
flips.

\subsection{Alerting Decisions Worth Defending}
\label{sec:impl-alerting}

\textbf{Kafka lag comes from Spark's own query progress, not a lag
exporter.} Structured Streaming does not commit consumer-group offsets
--- it tracks them in its checkpoint --- so \code{kafka-consumer-groups.sh}
and every off-the-shelf exporter report \emph{nothing} for these
consumers. The figure exists only in the progress event's source metrics.

\textbf{Batch metrics are exported by the API from PostgreSQL, not
pushed.} Airflow tasks are short-lived processes; Prometheus pulls, so by
scrape time the task has exited. A Pushgateway would work but is another
service holding job state indefinitely with no notion of staleness, and
is explicitly discouraged for this pattern~\cite{prometheus_docs}. The
batch layer already persists its outcomes because the API and the report
need them, so re-exporting costs nothing and guarantees the metric and
the report can never disagree.

Ten alert rules are defined (Appendix~\ref{app:alerts}), each with a
\code{for:} duration, a severity and a runbook entry.
Section~\ref{sec:alerting-results} reports which have been verified
firing and delivered.
